% GENERATED FILE. Edit canonical lessons, then run npm run generate:books.
% canonical-source-hash: fnv1a64:7990571d9c8157a7
% canonical-lessons: SA-C92-nimesah, SA-C92-purvam, SA-C92-pascat, SA-C92-adhuna, SA-C92-sarvada

\chapter{Moment, Before, After, Now, Always}
\label{ch:sa-moment-before-after-now-always}
\hlchaptermodality{92}

\begin{chapteropening}
\textbf{By the end of this chapter:} \emph{I can say a moment, before, after, now and always.}

Complete the last lesson of chapter 92: i can say a moment, before, after, now and always.
\end{chapteropening}

\section[nimeṣaḥ]{\sk{निमेषः} (nimeṣaḥ) — a moment, a blink}

\label{lesson:SA-C92-nimesah}



\begin{quote}
Before the new one: say the Sanskrit for late winter, then the Sanskrit for autumn.
\end{quote}

\subsection*{You'll want to know: \sk{निमेषः}}
\textbf{\sk{निमेषः}} — \emph{nimeṣaḥ} — \textquotedblleft{}a moment, a blink\textquotedblright{}.

\begin{grammarlens}[title={What you've built}]
One more word for this chapter.
\end{grammarlens}

\subsection*{Your turn}
\begin{itemize}
\raggedright
  \item \emph{Say it:} \emph{nimeṣaḥ}
  \item \emph{Say it:} \emph{nimeṣaḥ}, once more
  \item \emph{Say it:} say \emph{nimeṣaḥ}
  \item \emph{Recall:} say the Sanskrit for late winter, then the Sanskrit for autumn, then say \emph{nimeṣaḥ} again
\end{itemize}

\begin{tcolorbox}[breakable,colback=teal!4,colframe=teal!35!black,title={Before you move on}]
What does \sk{निमेषः} mean? (\textquotedblleft{}A moment, a blink\textquotedblright{}.) Say it once more.
\end{tcolorbox}
\section[pūrvam]{\sk{पूर्वम्} (pūrvam) — before, earlier}

\label{lesson:SA-C92-purvam}



\begin{quote}
Before the new one: say the Sanskrit for autumn, then the Sanskrit for a moment.
\end{quote}

\subsection*{You'll want to know: \sk{पूर्वम्}}
\textbf{\sk{पूर्वम्}} — \emph{pūrvam} — \textquotedblleft{}before, earlier\textquotedblright{}.

\begin{grammarlens}[title={What you've built}]
One more word for this chapter.
\end{grammarlens}

\subsection*{Your turn}
\begin{itemize}
\raggedright
  \item \emph{Say it:} \emph{pūrvam}
  \item \emph{Say it:} \emph{pūrvam}, once more
  \item \emph{Say it:} say \emph{pūrvam}
  \item \emph{Recall:} say the Sanskrit for autumn, then the Sanskrit for a moment, then say \emph{pūrvam} again
\end{itemize}

\begin{tcolorbox}[breakable,colback=teal!4,colframe=teal!35!black,title={Before you move on}]
What does \sk{पूर्वम्} mean? (\textquotedblleft{}Before, earlier\textquotedblright{}.) Say it once more.
\end{tcolorbox}
\section[paścāt]{\sk{पश्चात्} (paścāt) — after, later}

\label{lesson:SA-C92-pascat}



\begin{quote}
Before the new one: say the Sanskrit for a moment, then the Sanskrit for before.
\end{quote}

\subsection*{You'll want to know: \sk{पश्चात्}}
\textbf{\sk{पश्चात्}} — \emph{paścāt} — \textquotedblleft{}after, later\textquotedblright{}.

\begin{grammarlens}[title={What you've built}]
One more word for this chapter.
\end{grammarlens}

\subsection*{Your turn}
\begin{itemize}
\raggedright
  \item \emph{Say it:} \emph{paścāt}
  \item \emph{Say it:} \emph{paścāt}, once more
  \item \emph{Say it:} say \emph{paścāt}
  \item \emph{Recall:} say the Sanskrit for a moment, then the Sanskrit for before, then say \emph{paścāt} again
\end{itemize}

\begin{tcolorbox}[breakable,colback=teal!4,colframe=teal!35!black,title={Before you move on}]
What does \sk{पश्चात्} mean? (\textquotedblleft{}After, later\textquotedblright{}.) Say it once more.
\end{tcolorbox}
\section[adhunā]{\sk{अधुना} (adhunā) — now}

\label{lesson:SA-C92-adhuna}



\begin{quote}
Before the new one: say the Sanskrit for before, then the Sanskrit for after.
\end{quote}

\subsection*{You'll want to know: \sk{अधुना}}
\textbf{\sk{अधुना}} — \emph{adhunā} — \textquotedblleft{}now\textquotedblright{}.

\begin{grammarlens}[title={What you've built}]
One more word for this chapter.
\end{grammarlens}

\subsection*{Your turn}
\begin{itemize}
\raggedright
  \item \emph{Say it:} \emph{adhunā}
  \item \emph{Say it:} \emph{adhunā}, once more
  \item \emph{Say it:} say \emph{adhunā}
  \item \emph{Recall:} say the Sanskrit for before, then the Sanskrit for after, then say \emph{adhunā} again
\end{itemize}

\begin{tcolorbox}[breakable,colback=teal!4,colframe=teal!35!black,title={Before you move on}]
What does \sk{अधुना} mean? (\textquotedblleft{}Now\textquotedblright{}.) Say it once more.
\end{tcolorbox}
\section[sarvadā]{\sk{सर्वदा} (sarvadā) — always}

\label{lesson:SA-C92-sarvada}



\begin{quote}
Before the new one: say the Sanskrit for after, then the Sanskrit for now.

Say \textquotedblleft{}always.\textquotedblright{} (\emph{Sadā}.) This lesson's word says it another way.
\end{quote}

\subsection*{You'll want to know: \sk{सर्वदा}}
\textbf{\sk{सर्वदा}} — \emph{sarvadā} — \textquotedblleft{}always\textquotedblright{}.

\begin{grammarlens}[title={What you've built}]
One more word for this chapter.
\end{grammarlens}

\subsection*{Your turn}
\begin{itemize}
\raggedright
  \item \emph{Say it:} \emph{sarvadā}
  \item \emph{Say it:} \emph{sarvadā}, once more
  \item \emph{Say it:} say \emph{sarvadā}
  \item \emph{Recall:} say the Sanskrit for after, then the Sanskrit for now, then say \emph{sarvadā} again
\end{itemize}

\begin{tcolorbox}[breakable,colback=teal!4,colframe=teal!35!black,title={Before you move on}]
What does \sk{सर्वदा} mean? (\textquotedblleft{}Always\textquotedblright{}.) Say it once more.
\end{tcolorbox}
